\documentclass[12pt,a4paper]{article}

\usepackage[a4paper,margin=1in]{geometry}
\usepackage[T1]{fontenc}
\usepackage[utf8]{inputenc}
\usepackage{enumitem}
\setlength{\emergencystretch}{3em}

\title{\textbf{Supervisor Meeting Minutes}\\Meeting 6}
\author{Present: Matthew Chalmers, Jonas Leipner}
\date{Date: 3 November 2026}

\begin{document}

\maketitle

\section{Achieved in Previous Week}

\noindent\textit{Previous week: 27 October--3 November 2026}
\begin{itemize}[leftmargin=*]
    \item Updated the supervisor on the current state of the project repository,
    its documentation, and its planning artefacts.
    \item Continued refining the user stories and the possible direction of the
    project.
\end{itemize}

\section{Issues Discussed \& Solutions Agreed}

\begin{itemize}[leftmargin=*]
    \item \textbf{Repository and submission:} The repository structure and
    current project work were discussed. The final submission is expected to
    include a source-code ZIP, with the exact submission and repository-access
    requirements to be confirmed from the module guidance.
    \item \textbf{User stories:} The current user stories were reviewed as the
    basis for the game, research direction, and future implementation work.
    \item \textbf{Social aspect:} The previously implied multi-user and social
    dimension of the project was discussed explicitly. This includes learning
    from other users' towers, modules, configurations, and outcomes rather
    than necessarily making the game a direct player-versus-player game.
    Continuing with this social aspect was formally approved as part of the
    project direction.
    \item \textbf{Sources of shared content:} The discussion identified open
    questions about where modules or towers originate and how users can
    discover which options may be useful for them. Possible future mechanisms
    include shared modules, comparisons of successful configurations, and
    recommendations.
    \item \textbf{Sandbox communication frequency:} The next technical
    investigation is to determine when the \texttt{dart\_eval} sandbox should
    receive information from the game. Options include once per game loop,
    once whenever an enemy is present, once per relevant combat event, or only
    when an enemy enters a tower's range.
    \item \textbf{Sandbox input scope:} The experiment should compare passing
    information about one enemy, all enemies within a tower's range, or all
    enemies currently in the game. The choice should balance the information
    required by the dynamic logic against performance and implementation
    feasibility.
    \item \textbf{Evaluation value:} Measuring the communication frequency and
    input scope can provide evidence for the dissertation, including a
    comparison of the technical cost and practicality of alternative designs.
\end{itemize}

\section{Objectives for This Week}

\noindent\textit{Upcoming week: 3--10 November 2026}
\begin{itemize}[leftmargin=*]
    \item Create an initial experiment for passing combat information from the
    Flame game to the \texttt{dart\_eval} sandbox.
    \item Compare candidate invocation frequencies, including per-game-loop,
    per-event, and range-triggered execution.
    \item Compare candidate input scopes, including individual enemies,
    enemies within range, and all active enemies.
    \item Record measurable observations about feasibility, execution cost,
    and the information available to the dynamic damage calculation.
\end{itemize}

\section*{\small AI-use Disclaimer}
\small The meeting content and source notes for these minutes are genuine.
Artificial intelligence was used only to organise and summarise the notes
and transcript into a more digestible format; it was not used to invent the
meeting content or decisions.

\end{document}
